\chapter{Conclusion and Future Work}
\label{chap:conclusion_futurework}

This work addressed the challenge of automating physical design decisions for \ac{IVM} in PostgreSQL. While \ac{IVM} offers a compelling trade-off between read latency and data freshness, its write-side overhead is non-trivial and highly workload-dependent. We proposed, implemented, and evaluated a Machine Learning-based advisor capable of navigating this trade-off to recommend the optimal view strategy---No View, Materialized View, or Incremental View---for arbitrary SQL workloads.

\section{Conclusion}
\label{sec:conclusion}

The primary contributions of this work were the design of the advisor framework, the development of a novel write-cost estimation model that addresses the specific challenges of incremental maintenance, and the empirical evaluation of their zero-shot capabilities.

We successfully designed a view advisor that treats view selection as a cost-benefit optimization problem. By integrating read-cost estimates with our novel write-cost predictions, the advisor constructs a holistic view of the workload's performance profile. Crucially, our research highlights that accurate IVM advising requires observing internal trigger logic which is not exposed by standard \texttt{EXPLAIN} commands. Consequently, we designed the advisor to operate as a "What-If" analysis tool, leveraging a staging or sampling phase to capture the necessary trigger plans. The system demonstrates that learned models can effectively trade off read acceleration against write penalties, provided this planning phase is accommodated.

A core technical contribution of this thesis is the adaptation of the \ac{T3} cost modeling approach for \ac{IVM} maintenance. We demonstrated that standard tuple-centric prediction targets fail for maintenance workloads due to the presence of fixed overheads and non-linear scaling. To overcome this, we introduced a model that predicts per-pipeline \textit{View Refresh Time}, augmented by a composite feature engineering strategy. This strategy combines \textit{Delta Features} (from trigger execution plans) with \textit{Static Context Features} (from view definitions). This architecture allows the model to capture the interaction between update volume and view size, an important factor for zero-shot transfer.

The evaluation on the held-out IMDB benchmark demonstrated the practical utility of the approach despite the challenges of zero-shot transfer.
\begin{itemize}
    \item \textbf{Workload Acceleration:} The ultimate validation of the system is the aggregate impact on query processing. The advisor's recommendations resulted in a \textbf{$3.86\times$ speedup} for the total workload runtime compared to a baseline of fully materialized views. This confirms that the advisor successfully identifies and filters out scenarios where the cost of \ac{IMV} maintenance outweighs the read benefits.
    \item \textbf{Robust Decision Making:} While strict classification accuracy was limited to $41.67\%$ due to estimation noise in "close call" scenarios, the advisor achieved a \textbf{$91.67\%$ relaxed accuracy} within a 25\% tolerance threshold. This indicates that the system is highly effective at risk avoidance, consistently selecting strategies that are either optimal or competitively close to optimal.
    \item \textbf{Generalization and Calibration:} The write-cost model achieved a median Q-Error of $3.66\times$ on the unseen IMDB workload. The consistency between the median and $95^{th}$ percentile error ($3.72\times$) suggests that while the model successfully ranks relative costs, it suffers from a systematic scaling bias due to lack of input features on hardware overhead. This finding points toward the necessity of a "few-shot" calibration step in future work to align the model's absolute predictions with the target environment.
\end{itemize}

\section{Future Work}
\label{sec:future_work}

While this work establishes a strong foundation for learned \ac{IVM} advisors, several avenues for future research remain to enhance the system's accuracy and practicality.

\paragraph{Trigger Plan Estimation without Execution}
Currently, extracting features for the write-cost model requires executing \ac{DML} statements on a staging environment to capture the \ac{IVM} trigger plans. This introduces significant overhead. Future work should investigate modifying the PostgreSQL engine or the \texttt{pg\_ivm} extension to expose "Estimated Trigger Plans" via the standard \texttt{EXPLAIN} command. This would allow the advisor to operate in a purely hypothetical mode, drastically reducing the time required to generate recommendations.

\paragraph{Advanced Plan Representation}
The current model uses a flat, aggregate feature representation of the query plan. While effective for estimating the order of magnitude, the low Spearman rank correlation (0.08) indicates a limitation in fine-grained ranking. Future research could explore Graph Neural Networks (GNNs) or Tree-LSTM architectures to encode the full structure of the trigger plans. These deep learning approaches may better capture subtle operator interactions that tree-based models miss, potentially improving the ranking accuracy in zero-shot scenarios.

\paragraph{Concurrency and Lock Contention}
The current cost model assumes serial execution and focuses on resource consumption (CPU/IO). However, in high-throughput transactional workloads, the primary bottleneck for \ac{IMV} is often lock contention on the view or the delta tables. Developing a "Concurrency-Aware" cost model that predicts lock wait times based on write frequency and transaction isolation levels would be a critical step toward deploying \ac{IMV} advisors in operational OLTP environments.

\paragraph{Automatic Calibration}
Our evaluation revealed a systematic bias when transferring models between disparate hardware environments (e.g., "Tournament" training data vs. IMDB test environment). Future iterations of the advisor could implement an online learning module that runs a small sample of "probe" queries on the target database (Few-Shot Learning). By comparing the predicted vs. actual costs of these probes, the system could automatically calibrate its scaling factors, reducing the systematic prediction error and improving decision accuracy without full retraining.
